% GENERATED FILE. Edit canonical lessons, then run npm run generate:books.
% canonical-source-hash: fnv1a64:8b83a3dd5229c0f6
% canonical-lessons: HI-C329-chirakna, HI-C329-bhigona, HI-C329-uchalna, HI-C329-jhumna, HI-C329-masti

\chapter{A Festival of Colours}
\label{ch:hi-a-festival-of-colours}
\hlchaptermodality{329}

\begin{chapteropening}
\textbf{By the end of this chapter:} \emph{I can tell how a spring festival was played.}

Complete the last lesson of chapter 329: I can tell how a spring festival was played.
\end{chapteropening}

\section[\texorpdfstring{\hi{छिड़कना} (chiṛaknā)}{छिड़कना (chiṛaknā)}]{\hi{छिड़कना} (chiṛaknā) — to sprinkle}

\label{lesson:HI-C329-chirakna}



\begin{quote}
Before the new one: say the Hindi for a firecracker, then the Hindi for fireworks.
\end{quote}

\subsection*{You'll want to know: \hi{छिड़कना}}
\textbf{\hi{छिड़कना}} — \emph{chiṛaknā} — \textquotedblleft{}to sprinkle\textquotedblright{}.

\begin{quote}
\hi{दादी} \hi{ने} \hi{आँगन} \hi{में} \hi{पानी} \hi{छिड़का।}
\end{quote}

\begin{quote}
\emph{dādī ne āṁgan meṁ pānī chiṛkā.}
\end{quote}

\begin{quote}
Grandmother sprinkled water in the courtyard.
\end{quote}

Past, with \hi{ने} (he or she …): \hi{उसने} \hi{पानी} \hi{छिड़का}, \emph{usne pānī chiṛkā}.

\begin{cousinweb}
is inherited from \emph{*chiṭ}, a form CDIAL reconstructs for sudden splashes and scattering (CDIAL 5035). You \textbf{\hi{छिड़कना}} water, salt or colour.
\end{cousinweb}

\begin{grammarlens}[title={What you've built}]
One more word for telling what happened.
\end{grammarlens}

\subsection*{Your turn}
\begin{itemize}
\raggedright
  \item \emph{Say it:} \emph{chiṛaknā}
  \item \emph{Say it:} \emph{chiṛaknā}, once more
  \item \emph{Say it:} \emph{dādī ne āṁgan meṁ pānī chiṛkā.} — Grandmother sprinkled water in the courtyard.
  \item \emph{Recall:} say the Hindi for a firecracker, then the Hindi for fireworks, then say \emph{chiṛaknā} again
\end{itemize}

\begin{tcolorbox}[breakable,colback=teal!4,colframe=teal!35!black,title={Before you move on}]
What does \hi{छिड़कना} mean? (\textquotedblleft{}To sprinkle\textquotedblright{}.) Say it once more.
\end{tcolorbox}
\section[\texorpdfstring{\hi{भिगोना} (bhigonā)}{भिगोना (bhigonā)}]{\hi{भिगोना} (bhigonā) — to drench; to soak}

\label{lesson:HI-C329-bhigona}



\begin{quote}
Before the new one: say the Hindi for fireworks, then the Hindi for to sprinkle.
\end{quote}

\subsection*{You'll want to know: \hi{भिगोना}}
\textbf{\hi{भिगोना}} — \emph{bhigonā} — \textquotedblleft{}to drench; to soak\textquotedblright{}.

\begin{quote}
\hi{बच्चों} \hi{ने} \hi{बाल्टी} \hi{के} \hi{पानी} \hi{से} \hi{अरुण} \hi{को} \hi{भिगोया।}
\end{quote}

\begin{quote}
\emph{baccoṁ ne bālṭī ke pānī se Aruṇ ko bhigoyā.}
\end{quote}

\begin{quote}
The children drenched Arun with a bucket of water.
\end{quote}

Past, with \hi{ने} (he or she …): \hi{उसने} \hi{कपड़ा} \hi{भिगोया}, \emph{usne kapṛā bhigoyā}.

\begin{cousinweb}
\textbf{\hi{भिगोना}} (\emph{bhigonā}) is the causative of \emph{bhīgnā}, to get wet, which CDIAL 9500 sets under \emph{*bhiyagna}: to make someone or something wet through.
\end{cousinweb}

\begin{grammarlens}[title={What you've built}]
One more word for telling what happened.
\end{grammarlens}

\subsection*{Your turn}
\begin{itemize}
\raggedright
  \item \emph{Say it:} \emph{bhigonā}
  \item \emph{Say it:} \emph{bhigonā}, once more
  \item \emph{Say it:} \emph{baccoṁ ne bālṭī ke pānī se Aruṇ ko bhigoyā.} — The children drenched Arun with a bucket of water.
  \item \emph{Recall:} say the Hindi for fireworks, then the Hindi for to sprinkle, then say \emph{bhigonā} again
\end{itemize}

\begin{tcolorbox}[breakable,colback=teal!4,colframe=teal!35!black,title={Before you move on}]
What does \hi{भिगोना} mean? (\textquotedblleft{}To drench; to soak\textquotedblright{}.) Say it once more.
\end{tcolorbox}
\section[\texorpdfstring{\hi{उछालना} (uchālnā)}{उछालना (uchālnā)}]{\hi{उछालना} (uchālnā) — to toss up; to throw into the air}

\label{lesson:HI-C329-uchalna}



\begin{quote}
Before the new one: say the Hindi for to sprinkle, then the Hindi for to drench.
\end{quote}

\subsection*{You'll want to know: \hi{उछालना}}
\textbf{\hi{उछालना}} — \emph{uchālnā} — \textquotedblleft{}to toss up; to throw into the air\textquotedblright{}.

\begin{quote}
\hi{लोगों} \hi{ने} \hi{हवा} \hi{में} \hi{रंग} \hi{उछाला।}
\end{quote}

\begin{quote}
\emph{logoṁ ne havā meṁ raṅg uchālā.}
\end{quote}

\begin{quote}
People tossed colour into the air.
\end{quote}

Past, with \hi{ने} (he or she …): \hi{उसने} \hi{गेंद} \hi{उछाली}, \emph{usne gẽd uchālī}.

\begin{cousinweb}
is inherited from \emph{*ut-śālayati}, a form CDIAL reconstructs (CDIAL 1848), with the sense of throwing up. You \textbf{\hi{उछालना}} a ball, a coin or colour.
\end{cousinweb}

\begin{grammarlens}[title={What you've built}]
One more word for telling what happened.
\end{grammarlens}

\subsection*{Your turn}
\begin{itemize}
\raggedright
  \item \emph{Say it:} \emph{uchālnā}
  \item \emph{Say it:} \emph{uchālnā}, once more
  \item \emph{Say it:} \emph{logoṁ ne havā meṁ raṅg uchālā.} — People tossed colour into the air.
  \item \emph{Recall:} say the Hindi for to sprinkle, then the Hindi for to drench, then say \emph{uchālnā} again
\end{itemize}

\begin{tcolorbox}[breakable,colback=teal!4,colframe=teal!35!black,title={Before you move on}]
What does \hi{उछालना} mean? (\textquotedblleft{}To toss up; to throw into the air\textquotedblright{}.) Say it once more.
\end{tcolorbox}
\section[\texorpdfstring{\hi{झूमना} (jhūmnā)}{झूमना (jhūmnā)}]{\hi{झूमना} (jhūmnā) — to sway (with joy or music)}

\label{lesson:HI-C329-jhumna}



\begin{quote}
Before the new one: say the Hindi for to drench, then the Hindi for to toss up.
\end{quote}

\subsection*{You'll want to know: \hi{झूमना}}
\textbf{\hi{झूमना}} — \emph{jhūmnā} — \textquotedblleft{}to sway (with joy or music)\textquotedblright{}.

\begin{quote}
\hi{ढोल} \hi{बज} \hi{रहा} \hi{था} \hi{और} \hi{सब} \hi{लोग} \hi{झूम} \hi{रहे} \hi{थे।}
\end{quote}

\begin{quote}
\emph{ḍhol baj rahā thā aur sab log jhūm rahe the.}
\end{quote}

\begin{quote}
The drum was playing and everyone was swaying.
\end{quote}

Past (he …): \hi{वह} \hi{झूमा}, \emph{vah jhūmā}.

\begin{cousinweb}
\textbf{\hi{झूमना}} (\emph{jhūmnā}) is set by CDIAL under \emph{*kṣumbhati}, \textquotedblleft{}shakes\textquotedblright{} (CDIAL 3726), with a query of its own: the link is not certain. It is a happy, rocking movement: trees in the wind and people at a festival \textbf{\hi{झूमते} \hi{हैं}}.
\end{cousinweb}

\begin{grammarlens}[title={What you've built}]
One more word for telling what happened.
\end{grammarlens}

\subsection*{Your turn}
\begin{itemize}
\raggedright
  \item \emph{Say it:} \emph{jhūmnā}
  \item \emph{Say it:} \emph{jhūmnā}, once more
  \item \emph{Say it:} \emph{ḍhol baj rahā thā aur sab log jhūm rahe the.} — The drum was playing and everyone was swaying.
  \item \emph{Recall:} say the Hindi for to drench, then the Hindi for to toss up, then say \emph{jhūmnā} again
\end{itemize}

\begin{tcolorbox}[breakable,colback=teal!4,colframe=teal!35!black,title={Before you move on}]
What does \hi{झूमना} mean? (\textquotedblleft{}To sway (with joy or music)\textquotedblright{}.) Say it once more.
\end{tcolorbox}
\section[\texorpdfstring{\hi{मस्ती} (mastī)}{मस्ती (mastī)}]{\hi{मस्ती} (mastī) — merriment; high spirits}

\label{lesson:HI-C329-masti}



\begin{quote}
Before the new one: say the Hindi for to toss up, then the Hindi for to sway (with joy or music).
\end{quote}

\subsection*{You'll want to know: \hi{मस्ती}}
\textbf{\hi{मस्ती}} — \emph{mastī} — \textquotedblleft{}merriment; high spirits\textquotedblright{}.

\begin{quote}
\hi{त्योहार} \hi{के} \hi{दिन} \hi{बच्चे} \hi{सुबह} \hi{से} \hi{मस्ती} \hi{में} \hi{थे।}
\end{quote}

\begin{quote}
\emph{tyohār ke din bacce subah se mastī meṁ the.}
\end{quote}

\begin{quote}
On the day of the festival the children were in high spirits from the morning on.
\end{quote}

Feminine.

\begin{cousinweb}
\textbf{\hi{मस्ती}} (\emph{mastī}) is \emph{mast}, carefree and merry, usually traced to Persian \emph{mast}, \textquotedblleft{}intoxicated\textquotedblright{}, with the noun ending \emph{-ī}: the joy of a day with no work.
\end{cousinweb}

\begin{grammarlens}[title={What you've built}]
One more word for telling what happened.
\end{grammarlens}

\subsection*{Your turn}
\begin{itemize}
\raggedright
  \item \emph{Say it:} \emph{mastī}
  \item \emph{Say it:} \emph{mastī}, once more
  \item \emph{Say it:} \emph{tyohār ke din bacce subah se mastī meṁ the.} — On the day of the festival the children were in high spirits from the morning on.
  \item \emph{Recall:} say the Hindi for to toss up, then the Hindi for to sway (with joy or music), then say \emph{mastī} again
\end{itemize}

\begin{tcolorbox}[breakable,colback=teal!4,colframe=teal!35!black,title={Before you move on}]
What does \hi{मस्ती} mean? (\textquotedblleft{}Merriment; high spirits\textquotedblright{}.) Say it once more.
\end{tcolorbox}
